% === Section 4: DC3S Runtime Architecture ===
\section{The \dcs{} Runtime Architecture}\label{sec:architecture}

The ORIUS kernel orchestrates five typed stages, each consuming the output of the
previous stage and producing a well-typed contract object.  We describe each stage,
its inputs and outputs, the typed contracts involved, and the algorithm that
implements it.

\subsection{Architecture Overview}

\begin{figure}[h]
\centering
\begin{tikzpicture}[
  node distance=0.7cm and 1.5cm,
  stage/.style={draw, rounded corners, minimum width=2.5cm, minimum height=0.8cm,
                fill=blue!8, align=center, font=\small\bfseries},
  data/.style={font=\scriptsize\itshape, text=gray!70!black},
  arr/.style={-{Stealth[length=3pt]}, thick}
]
\node[stage] (D) {Stage 1\\Detect};
\node[stage, right=of D] (Ca) {Stage 2\\Calibrate};
\node[stage, right=of Ca] (Co) {Stage 3\\Constrain};
\node[stage, right=of Co] (S) {Stage 4\\Shield};
\node[stage, right=of S] (Ce) {Stage 5\\Certify};

\draw[arr] (D) -- (Ca) node[data,midway,above]{$\wt, \mathrm{flags}$};
\draw[arr] (Ca) -- (Co) node[data,midway,above]{$\Ut$};
\draw[arr] (Co) -- (S) node[data,midway,above]{$\Aset$};
\draw[arr] (S) -- (Ce) node[data,midway,above]{$a_t^{\mathrm{rel}}$};

\node[below=0.3cm of D, font=\scriptsize, text=gray] {ObservationPacket};
\node[below=0.3cm of Ca, font=\scriptsize, text=gray] {ReliabilityAssessment};
\node[below=0.3cm of Co, font=\scriptsize, text=gray] {ObsConsistentStateSet};
\node[below=0.3cm of S, font=\scriptsize, text=gray] {SafeActionSet};
\node[below=0.3cm of Ce, font=\scriptsize, text=gray] {SafetyCertificate};
\end{tikzpicture}
\caption{The five stages of the \dcs{} kernel and their inter-stage typed contracts.}
\label{fig:dc3s_pipeline}
\end{figure}

\subsection{Stage 1: Detect}\label{sec:detect}

The detect stage assesses the trustworthiness of the incoming telemetry.

\paragraph{Input.}  An \texttt{ObservationPacket} containing the raw telemetry
$\ztt$, the parsed state dictionary, a UTC timestamp, and the observation
history length.

\paragraph{Algorithm.}  Three detector families run in parallel:
\begin{enumerate}[nosep]
  \item \textbf{Staleness detector}: Checks whether the telemetry timestamp is
        within a configurable freshness window.  If stale, a binary flag
        \texttt{is\_stale} is raised and $\wt$ is penalised.
  \item \textbf{Page--Hinkley drift detector}: Maintains a running mean of each
        sensor channel and detects upward or downward drift using a cumulative-sum
        test with configurable threshold $\lambda$ and tolerance
        $\delta$~\cite{page1954continuous,hinkley1971inference}.
  \item \textbf{History-based heuristics}: Computes the number of consecutive
        identical readings (freeze detection), the rate of out-of-order timestamps,
        and spike magnitude relative to a rolling window.
\end{enumerate}

\paragraph{Output.}  A \texttt{ReliabilityAssessment} containing:
\begin{itemize}[nosep]
  \item $\wt \in [0,1]$: the composite reliability weight, clamped to $[0,1]$.
  \item \texttt{flags}: a dictionary of boolean detector outputs.
  \item \texttt{drift\_flag}: a boolean indicating whether drift was detected.
  \item \texttt{drift\_meta}: metadata from the Page--Hinkley detector
        (cumulative sum, threshold, direction).
  \item \texttt{assumption\_tags}: the assumptions activated by this assessment
        (A2, A5, A6).
\end{itemize}

\subsection{Stage 2: Calibrate}\label{sec:calibrate}

The calibrate stage computes the observation-consistent state set $\Ut$.

\paragraph{Input.}  The parsed state, the uncertainty dictionary (conformal
prediction intervals per target per horizon), and the reliability weight $\wt$.

\paragraph{Observation Quality Estimation (OQE).}
The composite reliability weight is the product of three independent scores:
\begin{equation}\label{eq:oqe}
  w_t = w_t^{\mathrm{fresh}} \cdot w_t^{\mathrm{range}} \cdot w_t^{\mathrm{consist}},
\end{equation}
where:
\begin{itemize}[nosep]
  \item $w_t^{\mathrm{fresh}} = \exp(-\lambda_f \cdot s_t)$, with $\lambda_f$ the
        freshness decay rate and $s_t$ the staleness counter (seconds since last
        novel reading).
  \item $w_t^{\mathrm{range}} = \mathbf{1}[x_t \in [\underline{x}, \bar{x}]]$,
        the range-validity indicator.
  \item $w_t^{\mathrm{consist}}$ penalises readings that deviate from the
        physics-model prediction beyond a configurable tolerance.
\end{itemize}

\paragraph{Reliability-Adaptive Conformal (RAC) Inflation.}
The key formula that connects telemetry reliability to the uncertainty set width is
the \emph{RAC-Cert} inflation:
\begin{equation}\label{eq:rac_cert}
  \boxed{\;
  C_t^{\mathrm{RAC}}(\alpha)
  \;=\;
  C_t(\alpha) \cdot \bigl[1 + \kappa_r\,(1 - w_t) + \kappa_d \cdot d_t + \kappa_s \cdot s_t\bigr]
  \;}
\end{equation}
where $C_t(\alpha)$ is the base conformal interval at miscoverage level~$\alpha$,
$\kappa_r = 0.5$ is the reliability penalty, $\kappa_d = 0.3$ is the drift penalty,
$d_t$ is the drift indicator, and $\kappa_s = 0.05$ is the staleness scaling with
$s_t$ the staleness count in seconds.

In the implementation, the simplified inflation factor used is
\begin{equation}\label{eq:inflation}
  \alpha = \mathrm{clip}\!\bigl(1.0 + 0.7\,(1 - w_t) + 0.5\,s_{\mathrm{shift}},
  \; 1.0,\; 3.0\bigr),
\end{equation}
which bounds the inflation between~1$\times$ (fully trusted) and~3$\times$ (maximally
degraded).

\paragraph{Output.}  An \texttt{ObservationConsistentStateSet} with:
\begin{itemize}[nosep]
  \item \texttt{observed\_state}: the point estimate from the parsed telemetry.
  \item \texttt{lower\_bounds}, \texttt{upper\_bounds}: per-target dictionaries.
  \item \texttt{inflation}: the applied factor $\alpha$.
  \item \texttt{quantile}: the conformal quantile level.
  \item \texttt{reliability\_weight}: the $\wt$ that produced this set.
\end{itemize}

\subsection{Stage 3: Constrain}\label{sec:constrain}

The constrain stage derives the safe action set $\Aset$ from $\Ut$ and the domain
safety specification.

\paragraph{Input.}  The \texttt{ObservationConsistentStateSet} from Stage~2 and a
\texttt{SafetySpec} encoding domain-specific constraints and the fallback action.

\paragraph{Algorithm.}  For each constraint in the safety specification:
\begin{enumerate}[nosep]
  \item Evaluate the constraint predicate over the worst-case state within $\Ut$.
  \item If the constraint is violated under worst-case evaluation, mark it as
        binding.
  \item The safe action set is the intersection of all non-binding constraint
        regions.
\end{enumerate}

In the battery domain, constraints include SoC bounds ($\SOC_{\min} \le \SOC_t \le
\SOC_{\max}$) and power limits ($|P_t| \le \Pmax$).  In the AV domain, constraints
include minimum gap maintenance and ego-speed bounds.

\paragraph{Output.}  A \texttt{SafeActionSet} containing the admissible region
and the list of candidate actions with their tightened constraint evaluations.

\subsection{Stage 4: Shield}\label{sec:shield}

The shield stage intercepts any candidate action that falls outside $\Aset$ and
either repairs it or substitutes the domain-specific fallback.

\paragraph{Input.}  The candidate action $\astar = \pi(\ztt)$, the
\texttt{SafeActionSet}, and the \texttt{SafetySpec} (which provides the fallback).

\paragraph{Algorithm.}
\begin{algorithmic}[1]
\If{$\astar \in \Aset$}
  \State \Return $\astar$ \Comment{Action is safe; pass through}
\Else
  \State Attempt projection: $a' = \mathrm{proj}_{\Aset}(\astar)$
  \If{$a'$ satisfies all constraints}
    \State \Return $a'$ \Comment{Repaired action}
  \Else
    \State \Return $\asafe$ \Comment{Fallback action}
  \EndIf
\EndIf
\end{algorithmic}

\paragraph{Output.}  A \texttt{RepairDecision} containing the released action
$a_t^{\mathrm{rel}}$, whether it was passed through, repaired, or substituted,
and the delta $\lVert\astar - a_t^{\mathrm{rel}}\rVert$.

\subsection{Stage 5: Certify}\label{sec:certify}

The certify stage produces a \texttt{SafetyCertificate} that binds together all
previous stage outputs into an auditable record.

\paragraph{Output.}  A \texttt{SafetyCertificate} containing:
\begin{itemize}[nosep]
  \item A unique certificate UUID.
  \item The released action and its provenance (pass-through, repair, fallback).
  \item The observation packet, reliability assessment, and uncertainty set.
  \item The constraint evaluation results.
  \item A timestamp and domain identifier.
  \item A payload dictionary for domain-specific metadata.
\end{itemize}

The certificate is appended to an append-only governance log.  The
\texttt{ContractVerifier} (\cref{sec:contracts}) validates every certificate
against the contract stack before it is committed to the log.

\subsection{Adapter Interface: The \texttt{DomainInstantiation} Protocol}

Domain-specific behaviour is confined to an adapter that implements six typed
methods defined by the \texttt{DomainInstantiation} protocol:
\begin{enumerate}[nosep]
  \item \texttt{ingest\_telemetry(raw\_bytes) -> ObservationPacket}: parse raw
        telemetry into a typed observation packet.
  \item \texttt{compute\_oqe(history) -> ReliabilityAssessment}: evaluate the OQE
        triple \cref{eq:oqe} and produce the composite weight~$w_t$.
  \item \texttt{build\_uncertainty\_set(state, w\_t, conformal\_q) ->
        ObsConsistentStateSet}: inflate the base conformal interval using
        \cref{eq:rac_cert}.
  \item \texttt{tighten\_action\_set(U\_t, spec) -> SafeActionSet}: derive $\Aset$
        via worst-case constraint evaluation over~$\Ut$.
  \item \texttt{repair\_action(a\_star, A\_t, spec) -> RepairDecision}: project or
        substitute the candidate action.
  \item \texttt{emit\_certificate(stages) -> SafetyCertificate}: bind all stage
        outputs into an auditable certificate.
\end{enumerate}

The same \dcs{} kernel binary orchestrates both the battery adapter
(\texttt{src/orius/battery\_adapter.py}) and the AV adapter
(\texttt{src/orius/av\_adapter.py}).  No kernel code is modified when switching
domains.

\subsection{The DC3S Kernel Algorithm}

\Cref{alg:dc3s} formalises the five-stage pipeline as a single procedural algorithm.

\begin{algorithm}[h]
\caption{DC3S Universal Kernel: \texttt{execute\_universal\_step}}\label{alg:dc3s}
\begin{algorithmic}[1]
\Require Domain adapter $\mathcal{D}$, raw telemetry $\mathbf{z}_t^{\mathrm{raw}}$,
         candidate action $a_t^*$, safety spec $\mathcal{C}$,
         conformal quantile $q_\alpha$, history buffer $\mathcal{H}$
\Ensure  Released action $a_t^{\mathrm{rel}}$, SafetyCertificate $\mathcal{K}_t$
\Statex
\State \textbf{[Stage 1: Detect]}
\State $\mathrm{obs}_t \gets \mathcal{D}.\texttt{ingest\_telemetry}(\mathbf{z}_t^{\mathrm{raw}})$
\State $(w_t, \mathrm{flags}, d_t, s_t) \gets \mathcal{D}.\texttt{compute\_oqe}(\mathcal{H} \cup \{\mathrm{obs}_t\})$
\Statex
\State \textbf{[Stage 2: Calibrate]}
\State $C_t(\alpha) \gets \texttt{conformal\_quantile}(q_\alpha, \mathrm{obs}_t)$
\State $C_t^{\mathrm{RAC}} \gets C_t(\alpha) \cdot [1 + \kappa_r(1-w_t) + \kappa_d d_t + \kappa_s s_t]$
  \Comment{\cref{eq:rac_cert}}
\State $\Ut \gets [\mathrm{obs}_t.\hat{x} - C_t^{\mathrm{RAC}},\; \mathrm{obs}_t.\hat{x} + C_t^{\mathrm{RAC}}]$
\Statex
\State \textbf{[Stage 3: Constrain]}
\State $\Aset \gets \mathcal{D}.\texttt{tighten\_action\_set}(\Ut, \mathcal{C})$
\Statex
\State \textbf{[Stage 4: Shield]}
\If{$a_t^* \in \Aset$}
  \State $a_t^{\mathrm{rel}} \gets a_t^*$ \Comment{Pass-through}
\Else
  \State $a' \gets \mathrm{proj}_{\Aset}(a_t^*)$ \Comment{Repair attempt}
  \If{$a' \in \Aset$}
    \State $a_t^{\mathrm{rel}} \gets a'$
  \Else
    \State $a_t^{\mathrm{rel}} \gets a_{\mathrm{safe}}$ \Comment{Fallback}
  \EndIf
\EndIf
\Statex
\State \textbf{[Stage 5: Certify]}
\State $\mathcal{K}_t \gets \mathcal{D}.\texttt{emit\_certificate}(\mathrm{obs}_t, w_t, \Ut, \Aset, a_t^{\mathrm{rel}})$
\State Append $\mathcal{K}_t$ to governance log
\State \Return $(a_t^{\mathrm{rel}}, \mathcal{K}_t)$
\end{algorithmic}
\end{algorithm}

\subsection{Battery-Domain Constraint Tightening: FTIT}

In the battery domain, Stage~3 enforces the \emph{Fault-Tolerant Invariant Tube}
(FTIT), which propagates the uncertainty from Stage~2 through the energy dynamics:
\begin{align}
  \bar{E}_{t+k} &= \hat{E}_t + \delta_{\max} + k\,\eta^{+}\,P_{\max}^{+}\,\Delta t, \label{eq:ftit_upper}\\
  \underline{E}_{t+k} &= \hat{E}_t - \delta_{\max} - k\,P_{\max}^{-}/\eta^{-}\,\Delta t. \label{eq:ftit_lower}
\end{align}
The safe-action set for the battery is then:
\begin{equation}
  \Aset = \Bigl\{ (P^+, P^-) \;\Big|\;
    E_{\min} + q_t^{\mathrm{RAC}} \;\le\; \hat{E}_{t+1}(P^+, P^-) \;\le\;
    E_{\max} - q_t^{\mathrm{RAC}} \Bigr\},
\end{equation}
where $q_t^{\mathrm{RAC}}$ is the RAC-inflated safety buffer derived from
$C_t^{\mathrm{RAC}}$.

\subsection{Latency Profile}

\Cref{tab:dc3s_latency} reports the per-stage latency measured over 12{,}096
runtime steps (battery pipeline with 42 fault-injected scenarios).

\begin{table}[h]
\caption{DC3S per-stage latency profile (battery pipeline, 12{,}096 steps).}
\label{tab:dc3s_latency}
\centering\small
\begin{tabular}{lrrrr}
\toprule
\textbf{Stage} & \textbf{Mean (ms)} & \textbf{Median (ms)} & \textbf{P95 (ms)} & \textbf{Max (ms)} \\
\midrule
OQE (Detect)      & 0.026 & 0.024 & 0.030 & 0.210 \\
Drift detect      & 0.000 & 0.000 & 0.000 & 0.003 \\
Calibrate         & 0.012 & 0.011 & 0.015 & 0.180 \\
Repair (Shield)   & 0.002 & 0.002 & 0.003 & 0.090 \\
\midrule
\textbf{Full step}& \textbf{0.042} & \textbf{0.037} & \textbf{0.041} & \textbf{2.416} \\
\bottomrule
\end{tabular}
\end{table}

The mean full-step latency of 0.042\,ms is well within the 10-minute dispatch interval for
battery systems and the 100\,ms control loop for autonomous vehicles.  Even the P95
latency of 0.041\,ms leaves ample headroom for real-time deployment.

\subsection{Detailed Stage Walkthrough: Battery Example}

To make the five-stage pipeline concrete, we trace the execution of a single
battery dispatch step under a stale-sensor fault.

\paragraph{Step context.}  The true SoC is 450\,MWh.  The sensor has been frozen for
3 intervals, so the telemetry reports \texttt{soc\_mwh = 440} (10~MWh stale error).
The promoted battery witness policy proposes a 50~MW discharge.

\paragraph{Stage 1 (Detect).}
The staleness detector records 3 consecutive identical readings and raises
\texttt{is\_stale = true}.  The reliability weight drops to $\wt = 0.976$.  The
Page--Hinkley detector does not fire because the constant readings mimic a
zero-drift signal.  Output: $\wt = 0.976$, \texttt{flags = \{is\_stale: true\}}.

\paragraph{Stage 2 (Calibrate).}
The battery runtime calibrator emits the base interval $\Chat = [420, 480]$\,MWh
(half-width 30~MWh).  With $\wt = 0.976$ and $s_{\mathrm{shift}} = 0$:
\[
  \alpha = \mathrm{clip}\!\bigl(1.0 + 0.7 \times 0.024 + 0.5 \times 0,\; 1.0,\; 3.0\bigr
  ) = 1.017.
\]
The inflated interval is $\Crac = [419.5, 480.5]$\,MWh.  The true SoC of 450~MWh
falls within $\Crac$, satisfying T1.

\paragraph{Stage 3 (Constrain).}
The safety specification requires $\SOC_{\min} = 100$~MWh $\le \SOC \le$
$\SOC_{\max} = 900$~MWh.  Even under worst-case evaluation ($\SOC = 419.5$~MWh),
a 50~MW discharge for one interval (10~min) reduces SoC by $\sim$8.3~MWh to
411.2~MWh, still above $\SOC_{\min}$.  The action is admitted to $\Aset$.

\paragraph{Stage 4 (Shield).}
Since the candidate action (50~MW discharge) is in $\Aset$, the shield passes it
through without repair.  Output: \texttt{pass\_through}, delta $= 0$.

\paragraph{Stage 5 (Certify).}
A certificate is emitted with UUID, the released action (50~MW discharge), the
stale-sensor flag, $\wt = 0.976$, $\alpha = 1.017$, and the constraint evaluation
result.  The certificate passes all predicates CP1--CP6.

This walkthrough demonstrates that even moderate degradation (10~MWh stale error) is
handled gracefully: the interval inflates by 1.7\%, the constraint check passes, and
no intervention is needed.

\subsection{Detailed Stage Walkthrough: AV Example}

We trace a single AV step under a spike fault.

\paragraph{Step context.}  The ego vehicle is travelling at 12~m/s with a relative
gap of 45~m.  A spike fault injects a transient reading of $v_{\mathrm{ego}} = 35$
m/s (23~m/s above truth).

\paragraph{Stage 1 (Detect).}
The spike detector computes the z-score of 35~m/s against the rolling mean (11.5~m/s)
and standard deviation (2.3~m/s): $z = (35 - 11.5)/2.3 = 10.2$.  This far exceeds
the spike threshold ($z > 3$), so \texttt{spike\_detected = true} is raised and
$\wt$ drops to approximately 0.40.

\paragraph{Stage 2 (Calibrate).}
With $\wt = 0.40$ and $s_{\mathrm{shift}} = 0$:
\[
  \alpha = \mathrm{clip}\!\bigl(1.0 + 0.7 \times 0.60 + 0.5 \times 0,\; 1.0,\; 3.0\bigr
  ) = 1.42.
\]
The base conformal interval for 1\,s ego speed is $\Chat = [8.0, 16.0]$~m/s
(half-width 4.0~m/s centred on 12~m/s).  After inflation:
$\Crac = [6.3, 17.7]$~m/s.  The true speed of 12~m/s is contained.

\paragraph{Stage 3 (Constrain).}
Under worst-case evaluation ($v = 6.3$~m/s), the minimum following gap may be
violated because the lead vehicle could be closer than expected.  The tightened safe
action set restricts the ego acceleration to $\le 0.5$~m/s$^2$ (rather than the
controller's proposed 1.2~m/s$^2$).

\paragraph{Stage 4 (Shield).}
The candidate acceleration (1.2~m/s$^2$) exceeds the tightened bound (0.5~m/s$^2$).
The shield projects the action to the boundary: released action = 0.5~m/s$^2$.
Output: \texttt{repaired}, delta $= 0.7$~m/s$^2$.

\paragraph{Stage 5 (Certify).}
The certificate records the repair delta, the spike flag, $\wt = 0.40$,
$\alpha = 1.42$, and the constraint-tightened evaluation.

This walkthrough shows the full shielding pathway: a severe spike triggers a low
reliability weight, which inflates the interval by 42\%, which tightens the constraint
set, which triggers an action repair.

\subsection{Kernel Code Structure}

The ORIUS kernel codebase is organised as follows:

\begin{verbatim}
src/orius/
  universal_theory/
    kernel.py          # Five-stage orchestration (~350 lines)
    contracts.py       # Typed dataclasses + Protocol
    risk_bounds.py     # T10/T11 risk computations
  dc3s/
    drift.py           # Page-Hinkley drift detector
    runtime.py         # DC3S runtime loop
  battery_adapter.py   # Battery DomainInstantiation
  av_adapter.py        # AV DomainInstantiation
\end{verbatim}

The total kernel (excluding adapters) is under 600~lines of Python.  The two adapters
add approximately 150--200~lines each.  This compact implementation supports the claim
that the framework is lightweight and auditable.
